% =============================================================================
%  sec_feasibility.tex -- Section: Feasibility Assessment
%  BodyCheck -- AI-Powered Medical Triage & 3D Symptom Checker
%  Vietnamese Student HackAIthon 2026 -- Bang B (Challenger)
%  Team Clavicular | HCMUS APCS/CLC
%
%  Pure content module -- include via \input{sec_feasibility}.
%  All \definecolor declarations live in main.tex preamble.
%
%  Diagram images go in:  hackAIthon_latex/figures/
%    arch_diagram.png    -- Mermaid source: PLACE_arch_diagram_HERE.txt
%    gantt_roadmap.png   -- Mermaid source: PLACE_gantt_roadmap_HERE.txt
% =============================================================================

\section{Feasibility Assessment}
\label{sec:feasibility}

This section addresses the five feasibility sub-criteria defined in the third rubric of the contest --- Feasibility:
legal data sourcing and qualified personnel; technical build and deployment
viability; infrastructure and operational cost estimation; security and legal
compliance under current Vietnamese law; and the post-competition GTM
roadmap. Evidence is drawn from the team's direct engineering experience and
the organiser-provided VNPT API ecosystem that will power the Round~2 build.

Table~\ref{tab:feasibility_matrix} gives a structured overview; each dimension
is examined in depth in the subsections that follow.

% ---------------------------------------------------------------------------
% Table 1 -- Five-Axis Feasibility Matrix
% ---------------------------------------------------------------------------
\begin{table}[htbp]
  \centering
  \caption{BodyCheck five-axis feasibility assessment aligned to the
    HackAIthon 2026 Bang~B Round~1 rubric (25 pts total).
    Risk ratings: \textcolor{clav-green}{\textbf{Low}} = controls in place;
    \textcolor{clav-amber}{\textbf{Medium}} = mitigations identified.}
  \label{tab:feasibility_matrix}
  \setlength{\tabcolsep}{5pt}
  \renewcommand{\arraystretch}{1.45}
  \small
  \begin{tabular}{@{} p{2.6cm} p{5.0cm} p{1.6cm} p{3.6cm} @{}}
    \toprule
    \textbf{Dimension} &
    \textbf{Key Evidence} &
    \textbf{Risk} &
    \textbf{Current Status} \\
    \midrule

    Legal Data \& HR &
    No PHI stored; VNPT APIs under organiser-provided access;
    4-person HCMUS CS student team &
    \textcolor{clav-green}{\textbf{Low}} &
    Cleared for Round~1 proposal \\

    Tech Build \& Deploy &
    Three-tier Flutter / Node.js / FastAPI stack;
    VNPT API integration path confirmed &
    \textcolor{clav-green}{\textbf{Low}} &
    Round~2 MVP deliverable within 8 days \\

    Infra \& Costs &
    VNPT APIs free in Round~2 (organiser-provisioned);
    single low-cost compute instance post-competition &
    \textcolor{clav-amber}{\textbf{Medium}} &
    Post-competition billing TBD\textsuperscript{\dag} \\

    Security \& Legal &
    TLS enforced; vault secrets; AI disclaimer embedded;
    general info tool classification at MVP stage &
    \textcolor{clav-amber}{\textbf{Medium}} &
    Clinical/legal review before public release \\

    Roadmap \& GTM &
    Three-phase post-competition plan; B2C + B2B2C channels;
    12-month market expansion roadmap &
    \textcolor{clav-green}{\textbf{Low}} &
    Detailed in \S\ref{sec:feasibility-gtm} \\

    \bottomrule
  \end{tabular}

  \smallskip\noindent
  {\footnotesize \textsuperscript{\dag}Post-competition API pricing and
  compute cost estimates will be confirmed once commercial VNPT tier
  information is available.}
\end{table}

% =============================================================================
\subsection{Legal Data Acquisition and Human Resources}
\label{sec:feasibility-hr}
% =============================================================================

\subsubsection*{Data Sourcing and Privacy Compliance}

BodyCheck is architected to avoid handling protected health information (PHI)
at the infrastructure level. Symptom descriptions entered by users are treated
as transient, request-scoped payloads: they are forwarded to the AI reasoning
layer, a structured clinical summary is generated, and the raw input is
discarded without being written to any persistent store or log file in
identifiable form.

Medical reference content is retrieved at inference time from authoritative
public-facing sources via API rather than ingested into a proprietary dataset.
In Round~2 of the competition, this retrieval layer will be implemented using
\textbf{VNPT Smartbot} (advanced LLM-backed question answering, API~4.2) and
\textbf{VNPT SmartReader} (document OCR and information extraction, API~5.1/5.2),
both provided by the organiser under the competition's API access agreement.
This eliminates all data-licensing risk for the competition phase: every API
call is governed by VNPT's enterprise terms and is explicitly sanctioned for
participating teams by the organiser.

The system therefore operates without requiring independent data-collection
agreements, proprietary patient datasets, or any form of model fine-tuning on
personally identifiable health records.

\subsubsection*{Team Composition and Skill Coverage}

The team consists of four full-time undergraduate students from the
Faculty of Information Technology, Ho Chi Minh City University of Science
(HCMUS). Each member specializes in a distinct engineering
domain:

\begin{itemize}[leftmargin=2em, itemsep=2pt, topsep=4pt]
  \item \textbf{AI Engineering} (\textit{Nguyen Dinh Thien Loc}):
        LLM orchestration, VNPT Smartbot and SmartVoice integration,
        FastAPI AI microservice design, and prompt engineering.
  \item \textbf{Frontend / UI} (\textit{Quach Thien Lac}):
        Flutter mobile/web client, Three.js 3D anatomical canvas
        (iframe wrapper), and VNPT SmartUX analytics embedding.
  \item \textbf{3D Asset Design} (\textit{Nguyen Quoc Nam}):
        Blender modelling, GLB export and optimisation, and clickable
        anatomical region mapping across three model files.
  \item \textbf{Backend / Gateway} (\textit{Tran Ton Minh Ky}):
        Node.js/Express API gateway, JSON schema validation middleware,
        VNPT eKYC integration, and error-handling controllers.
\end{itemize}

All four members are currently enrolled full-time students, satisfying the
eligibility criteria for Bang~B -- Challenger (up to five members per team,
including an optional faculty advisor). The team is self-sufficient at the
MVP scale required for Rounds~1 and~2. Post-competition scaling toward a
production-grade platform would additionally require a DevOps/SRE engineer
and a licensed medical content adviser.

% As a validation point, the team built and demonstrated a working first
% prototype of the core triage loop --- 3D body-part selection, symptom input,
% AI-generated clinical summary, and clinic geo-routing --- in a 72-hour sprint
% at a prior student hackathon, confirming that the stack is viable and that
% the team can deliver under time pressure.

% =============================================================================
\subsection{Technical Build and Deployment Feasibility}
\label{sec:feasibility-tech}
% =============================================================================

\subsubsection*{Technology Stack Viability}

BodyCheck employs a \emph{decoupled three-tier service architecture} that
separates public-facing API traffic from the AI orchestration layer, making
each tier independently testable, replaceable, and deployable:

\begin{itemize}[leftmargin=2em, itemsep=2pt, topsep=4pt]
  \item \textbf{Flutter Web/Mobile (Client Tier):}
        A single Dart codebase compiles to native mobile (ARM) and
        optimised web (Wasm/JS), so the same product ships as a PWA and
        as a native app without a separate codebase. The interactive 3D
        anatomical model is rendered by a Three.js scene embedded as an
        iframe; region-selection events travel back to Flutter via
        \texttt{postMessage}.
  \item \textbf{Node.js 20 LTS / Express 5 (Gateway Tier):}
        Handles CORS, JWT authentication, JSON Schema request validation,
        and reverse-proxying to the AI microservice. Node's event-loop
        model absorbs concurrent upstream API calls without thread contention.
  \item \textbf{Python 3.12 / FastAPI (AI Orchestration Tier):}
        ASGI-native, Pydantic-validated. FastAPI provides auto-generated
        OpenAPI docs, sub-millisecond request parsing, and clean async
        bindings for all VNPT client SDKs.
  \item \textbf{VNPT API Ecosystem (Round~2):}
        The organiser provides direct access to eight API families.
        BodyCheck's planned integration uses:
        SmartVoice STT (API~3.2) for voice symptom input;
        SmartVoice TTS (API~3.1) for reading responses back to users
        with accessibility needs;
        SmartVoice call-summary (API~3.3) for distilling multi-turn
        consultation transcripts into a structured summary;
        Smartbot rule-based triage (API~4.1) for fast, deterministic
        first-pass classification;
        Smartbot advanced LLM (API~4.2) for open-ended differential
        reasoning and follow-up chat;
        SmartReader OCR and extraction (API~5.1/5.2) for parsing
        user-uploaded medical documents, lab reports, and prescriptions;
        SmartUX analytics (API~7.1) for in-app interaction telemetry
        and friction-point identification;
        vnSocial trend listening (API~6.1/6.2) for monitoring
        health-topic sentiment on Vietnamese social media, surfacing
        epidemic signals and common symptom clusters;
        eKYC document OCR (API~1.1) for verified account creation in
        the B2B2C institutional channel; and
        SmartVision face detection (API~8.3) as an optional
        telemedicine extension that confirms user presence before a
        live consultation handoff.
        Using the full breadth of the organiser-provided API catalogue
        maximises the integration score in the Round~2 product
        completion rubric and demonstrates the depth of the team's
        engagement with the VNPT ecosystem.
\end{itemize}

Figure~\ref{fig:arch-diagram} shows the complete deployment topology,
including the VNPT API integration layer.

% ---------------------------------------------------------------------------
% Figure 1 -- Architecture Diagram (Mermaid export)
% ---------------------------------------------------------------------------
\begin{figure}[htbp]
  \centering
  \IfFileExists{figures/arch_diagram.png}{%
    \includegraphics[width=0.95\textwidth, keepaspectratio]%
      {figures/arch_diagram.png}%
  }{%
    \fbox{%
      \begin{minipage}{\dimexpr0.95\textwidth-2\fboxsep-2\fboxrule\relax}
        \centering\vspace{2cm}
        {\large\sffamily\color{clav-gray}
          [Figure~1: System Architecture Diagram -- placeholder]\\[6pt]}
        {\small\sffamily\color{clav-gray}
          Export \texttt{arch\_diagram.png} from Mermaid
          (source: \texttt{PLACE\_arch\_diagram\_HERE.txt})
          and save to the \texttt{figures/} folder.}
        \vspace{2cm}
      \end{minipage}%
    }%
  }
  \caption{BodyCheck three-tier deployment topology with VNPT API integration.
    Arrows denote synchronous HTTPS calls between components; shaded
    subgraphs are independently containerisable service tiers.
    The AI Orchestration Tier calls VNPT SmartVoice, Smartbot, and
    SmartReader in Round~2; post-competition builds may extend to
    additional partner APIs.}
  \label{fig:arch-diagram}
\end{figure}

\subsubsection*{Local Development and CI/CD Workflow}

All three tiers run concurrently on a single developer laptop: Uvicorn on
port~8016, Express on port~3000, and \texttt{flutter run -d chrome} pointing
at the local gateway. A new contributor can clone the repository and have
every service running in under 15~minutes.

For the Round~2 submission, each service is containerised with a minimal
\texttt{Dockerfile} and composed via \texttt{docker-compose.yml},
satisfying the competition's one-command setup requirement for product
completion. A lightweight GitHub Actions workflow runs automated smoke tests
against the \texttt{/api/diagnose}, \texttt{/api/chat}, and
\texttt{/api/clinics} endpoints before each push to the submission branch.

% =============================================================================
\subsection{Infrastructure and Operational Costs}
\label{sec:feasibility-costs}

\subsubsection*{Round~2 --- Competition Phase}

During Round~2 (26~June -- 3~July~2026), the organiser provides all VNPT API
access at no cost to the team. Infrastructure expenditure is therefore
limited to the compute instance required to run the three-tier service stack
--- a university lab server or a minimal cloud VM (2~vCPU, 2~GB RAM) is
sufficient for the prototype load expected during judging and live-demo
sessions.

\subsubsection*{Post-Competition Cost Model}

After the competition, the team must transition to a commercial or partner
API pricing arrangement. The total monthly operational expenditure scales
linearly with session volume $R$:

\begin{equation}
  \label{eq:cost-model}
  C_{\mathrm{ops}}(R) \;=\; C_{\mathrm{fixed}} \;+\;
  \sum_{i=1}^{n} c_i\,\alpha_i\,R
\end{equation}

where $C_{\mathrm{fixed}}$ covers compute, networking, and storage;
$c_i$ is the per-call unit price of VNPT service~$s_i$; and $\alpha_i$
is the number of calls per session. Because the AI token budget per session
is bounded by a fixed prompt template, this cost model is
\emph{strictly linear} in $R$ with no quadratic blowup.

The break-even monthly session volume $R^{*}$ under monetisation price $p$
satisfies:

\begin{equation}
  \label{eq:break-even}
  R^{*} \;=\; \frac{C_{\mathrm{fixed}}}
               {p \;-\; \displaystyle\sum_{i=1}^{n} c_i\,\alpha_i}
\end{equation}

\noindent\textit{Note on specific figures:} Commercial per-call pricing
for VNPT API tiers is not yet confirmed. Detailed cost estimates will
be inserted once the team establishes a commercial arrangement with VNPT
or an equivalent partner post-competition. The structural model above
remains valid regardless of the specific coefficient values.

% =============================================================================
\subsection{Security and Legal Compliance}
\label{sec:feasibility-security}
% =============================================================================

\subsubsection*{Data Protection Architecture}

BodyCheck treats all symptom inputs as potentially sensitive health information,
applying defence-in-depth controls at every tier:

\begin{itemize}[leftmargin=2em, itemsep=3pt, topsep=4pt]
  \item \textbf{Transport Security:}
        All client-to-gateway and inter-service communication uses TLS~1.3.
        Internal service calls within a private VPC stay off the public internet.
  \item \textbf{Secrets Management:}
        API credentials (VNPT keys, gateway tokens) are injected as
        environment variables from a secrets vault (Docker Secrets or
        equivalent) and never appear in source control. The repository
        includes only a \texttt{.env.example} template listing required
        variable names without exposing values.
  \item \textbf{Input Validation:}
        The Express gateway applies JSON Schema middleware before any payload
        reaches the AI tier, blocking prompt-injection and malformed-input
        vectors at the network boundary.
  \item \textbf{No Persistent PHI:}
        The FastAPI service holds session context only within the lifespan of
        a single HTTP request. Structured logs record only anonymised metadata
        (request ID, latency, model tag, HTTP status). No free-text symptom
        content appears in log output.
  \item \textbf{Rate Limiting:}
        Per-IP rate limits and API-key quotas at the gateway prevent
        runaway API quota consumption and abuse of the organiser-provided
        VNPT allocation.
\end{itemize}

\subsubsection*{Open-Source Licensing and Medical Liability}

The BodyCheck codebase is released under the MIT License, which explicitly
disclaims all warranties and limits liability. Because the system produces
health-adjacent guidance, additional framing is mandatory:

\begin{itemize}[leftmargin=2em, itemsep=3pt, topsep=4pt]
  \item A medical disclaimer is embedded in every AI response payload and
        displayed prominently in the Flutter UI, clearly stating that
        BodyCheck provides \emph{general informational support only} and is
        not a substitute for professional medical advice.
  \item Emergency escalation text (Vietnam emergency line: \textbf{115})
        is injected by the prompt template whenever the AI detects a
        high-severity symptom pattern.
  \item Pre-launch clinical and legal review by a licensed medical
        professional and by legal counsel specialising in Vietnamese
        health-sector regulations is required before any public-facing
        deployment beyond the competition context.
\end{itemize}

\subsubsection*{Regulatory Classification}

Under current Vietnamese law, the relevant authorities are the Ministry of
Information and Communications (MOIT) for digital platform registration
and the Ministry of Health (MOH) for services classified as medical devices
or telemedicine under Circular~30/2020/TT-BYT.

At MVP stage, BodyCheck is a \emph{general health information tool}, not a
telemedicine service or medical device, provided the non-diagnostic framing
and mandatory disclaimer are consistently maintained. This keeps the
competition-phase product outside medical device registration requirements.
Future versions targeting structured diagnostic conclusions would require
registration under Decree~98/2021/ND-CP and may require ISO~13485
quality management certification.

% =============================================================================
\subsection{Post-Competition Roadmap and GTM Strategy}
\label{sec:feasibility-gtm}

BodyCheck's market scope is intentionally \emph{global from day one}: the
entire product interface, AI prompt layer, and clinical-reference content
operate in English, making it accessible to users in most countries without
a localisation bottleneck. Vietnam serves as the primary launch market given
the team's local regulatory familiarity and existing network, with South-East
Asia and broader English-speaking markets as the natural second wave. We plan to eventually include Vietnamese language to the product to be to reach older demographics domestically, but this is a post-MVP feature rather than a launch requirement.

Beyond the primary consumer health use case, BodyCheck addresses a material
problem in \textbf{education}: medical and nursing students, first-aid
training programmes, and health-literacy curricula routinely struggle to
teach anatomical symptom recognition interactively. The same 3D body-region
interface and triage-reasoning engine that serves a self-diagnosing user
also serves a student learning clinical decision pathways --- creating a
dual-market opportunity from a single codebase. This educational channel
aligns directly with HackAIthon topic~1 (AI solutions supporting
Personalised Learning) in addition to the primary topic~4 (AI supporting
Healthcare Operations).

The post-competition roadmap is structured in three phases aligned with
the rubric's 12-month market expansion requirement.

\subsubsection*{Phase Definitions}

\begin{description}[leftmargin=1em, style=nextline, itemsep=8pt, topsep=4pt]

  \item[\textbf{Phase 1 --- HackAIthon Competition}]
    \textit{June -- July 2026 (Rounds 1--3).}\\
    Submit the Round~1 idea proposal. If selected for Round~2
    (announcement before 26~June), integrate necessary VNPT API services
  into the BodyCheck stack over the 8-day build window
    (26~June -- 3~July) with direct Mentor support. Demonstrate a
    live-running product at the Round~3 finals (14--15~July~2026 in
    Hanoi). Collect structured judge and mentor feedback to anchor the
    post-competition product direction.

  \item[\textbf{Phase 2 --- Hardening and Closed Beta}]
    \textit{August -- November 2026 (Months 1--4 post-competition).}\\
    Begin serious post-competition development in August once semester
    schedules allow. Stabilise and document all VNPT API integrations;
    containerise all services with a one-command
    \texttt{docker-compose} setup; implement TLS certificate
    provisioning and vault-based secret injection.
    In September--October, run a closed beta with 30--50 volunteer
    testers drawn from HCMUS student and faculty networks and online
    medical communities. Iterate on VNPT Smartbot prompt templates and
    SmartReader parsing pipelines based on real-user feedback. Conduct
    a first-pass clinical language and legal review.
    Target: a stable, security-hardened beta release by end of
    November~2026.

  \item[\textbf{Phase 3 --- Public Launch and Market Expansion}]
    \textit{December 2026 -- June 2027 (Months 5--12 post-competition).}\\
    Add persistent session storage, user account management, and a UX
    metrics observability stack (structured logging, distributed
    tracing). Launch the consumer app (B2C) as a PWA and on mobile
    stores with a freemium model in February~2027. In parallel,
    onboard the first education-sector partners (medical schools, online
    health-literacy platforms) under a white-label API licence in
    March~2027, followed by the broader B2B2C healthcare channel
    (telehealth operators, hospital intake portals) in April~2027.
    Pursue institutional funding via accelerator programmes
    (VNPT Innovation Hub, Google for Startups Vietnam, HCMUS
    technology-transfer office) starting January~2027.
    Initiate MOH regulatory consultation for the clinical-grade product
    tier in April~2027.
    Target: stable public launch \textbf{v1.0} by June~2027,
    approximately 12~months post-competition.

\end{description}

\subsubsection*{Go-To-Market Channel Strategy}

BodyCheck targets three complementary GTM channels:

\begin{enumerate}[leftmargin=2em, itemsep=4pt, topsep=4pt]
  \item \textbf{B2C --- Global Consumer Health App (English-first):}
        A free-tier PWA and native mobile app distributed worldwide.
        The English-first UI removes localisation barriers and makes
        the product immediately accessible across South-East Asia,
        South Asia, and other regions underserved by AI-powered
        medical triage tools. Revenue via freemium subscription.
  \item \textbf{B2B2C --- API Licensing to Healthcare Platforms:}
        White-label REST API for telehealth operators, hospital patient
        intake portals, and health insurance triage systems. Priced at
        a per-call rate plus a monthly platform fee. Primary targets
        are Vietnamese healthcare organisations; secondary targets are
        international digital-health platforms that need a
        multi-language AI triage layer.
  \item \textbf{Education Channel --- Medical Training Platforms:}
        The same 3D anatomical model and AI reasoning engine that
        supports self-triage also serves medical and nursing students,
        first-aid training programmes, and health-literacy curricula.
        This channel is licensed to universities and online learning
        platforms as a white-label simulation tool, at a flat
        institutional annual fee rather than per-session billing.
\end{enumerate}

Figure~\ref{fig:gantt-roadmap} shows the full 12-month roadmap anchored
to the competition milestone.

% ---------------------------------------------------------------------------
% Figure 2 -- Gantt Roadmap (Mermaid export)
% ---------------------------------------------------------------------------
\begin{figure}[htbp]
  \centering
  \IfFileExists{figures/gantt_roadmap.png}{%
    \includegraphics[width=\textwidth, keepaspectratio]%
      {figures/gantt_roadmap.png}%
  }{%
    \fbox{%
      \begin{minipage}{\dimexpr\textwidth-2\fboxsep-2\fboxrule\relax}
        \centering\vspace{2cm}
        {\large\sffamily\color{clav-gray}
          [Figure~2: 12-Month Post-Competition Roadmap -- placeholder]\\[6pt]}
        {\small\sffamily\color{clav-gray}
          Export \texttt{gantt\_roadmap.png} from Mermaid
          (source: \texttt{PLACE\_gantt\_roadmap\_HERE.txt})
          and save to the \texttt{figures/} folder.}
        \vspace{2cm}
      \end{minipage}%
    }%
  }
  \caption{BodyCheck 12-month development and GTM roadmap anchored at the
    HackAIthon 2026 finals (July 2026). Phase~1 covers the competition
    period; Phases~2 and~3 represent the post-competition commercialisation
    trajectory. Completed tasks appear in green; planned tasks in blue;
    diamonds mark major milestones.}
  \label{fig:gantt-roadmap}
\end{figure}

% ---------------------------------------------------------------------------
% Table 2 -- Risk Matrix
% ---------------------------------------------------------------------------
\begin{table}[htbp]
  \centering
  \caption{Principal risks and mitigation strategies for the post-competition
    product roadmap.}
  \label{tab:risk-matrix}
  \setlength{\tabcolsep}{6pt}
  \renewcommand{\arraystretch}{1.4}
  \small
  \begin{tabular}{@{} p{3.1cm} p{2.1cm} p{7.0cm} @{}}
    \toprule
    \textbf{Risk} & \textbf{Severity} & \textbf{Mitigation Strategy} \\
    \midrule

    VNPT commercial API pricing &
    \textcolor{clav-amber}{\textbf{Medium}} &
    Provider-agnostic routing layer in FastAPI allows hot-swapping to
    alternative AI backends (e.g., open-source LLM hosts, Gemini API)
    without touching application code above the service boundary. \\[4pt]

    Regulatory reclassification &
    \textcolor{clav-amber}{\textbf{Medium}} &
    Phase~3 launch targets the non-regulated consumer information category.
    Clinical-grade features are deferred behind a separate product tier
    pending MOH review. The mandatory disclaimer ensures compliance
    throughout. \\[4pt]

    Team continuity &
    \textcolor{clav-amber}{\textbf{Medium}} &
    Infrastructure-as-Code, clean API boundaries between tiers, and
    comprehensive technical documentation lower the barrier for new
    contributors and reduce single-point-of-failure risk. \\[4pt]

    User trust in AI medical tools &
    \textcolor{clav-amber}{\textbf{Medium}} &
    Transparent uncertainty communication, mandatory disclaimer, and
    triage-support framing (not diagnosis) are embedded in the prompt
    template and UI copy at every touchpoint. \\

    \bottomrule
  \end{tabular}
\end{table}

% =============================================================================
% End of Section: Feasibility Assessment
% =============================================================================
